\hypertarget{appendix-d8a7ddb675eb}{}\label{referencestart:appendix-d8a7ddb675eb}
\section{Jacobians, pushforwards, and pullbacks}\label{appendix-d8a7ddb675eb-jacobians-pushforwards-and-pullbacks}

\textbf{Curriculum depth:} Developed · \textbf{First seen:} Phase 9 · \textbf{Latest development:} Phase 9

\begin{allmightroledefinition}

\subsection{Definition and mechanism}\label{appendix-d8a7ddb675eb-definition-and-mechanism}

The Jacobian is the matrix of first derivatives of a vector-valued function. It maps a small input displacement to a first-order output displacement, a forward sensitivity often called a pushforward or Jacobian-vector product. Gradients transform in the reverse direction by the transpose Jacobian, the pullback used in reverse-mode differentiation. These operations differ because vectors and covectors have different transformation rules. Jacobians compose by matrix multiplication and their determinants describe local volume changes. A Jacobian is local: it does not replace a nonlinear transformation globally. At singular points it can lose rank or become ill-conditioned, affecting inverse sensitivity, uncertainty propagation, and optimizer steps. State coordinate ordering and column-vector conventions before interpreting signs or multiplication order.

\end{allmightroledefinition}

\subsection{Prerequisites}\label{appendix-d8a7ddb675eb-prerequisites}

\begin{itemize}
\tightlist
\item
  \hyperref[appendix-7bc181935c42]{Derivative as a local linear map}
\item
  \hyperref[appendix-ab31ef447dbf]{Matrix multiplication and composition order}
\item
  \href{../02-Mathematics-for-Visual-Computing/main.pdf\#nameddest=appendix-4815cd87ab8b}{Probability density and change of measure (Mathematics for Visual Computing, Reference notes: Probability density and change of measure)}
\end{itemize}

\subsection{Used by}\label{appendix-d8a7ddb675eb-used-by}

\begin{itemize}
\tightlist
\item
  \href{../04-Light-Transport-and-Image-Synthesis/main.pdf\#nameddest=appendix-40df0855b4a7}{Automatic differentiation and backpropagation (Light Transport and Image Synthesis, Reference notes: Automatic differentiation and backpropagation)}
\item
  \hyperref[appendix-6ac071f8b50c]{Hessians and local curvature}
\item
  \hyperref[appendix-2c2df8345991]{Nonlinear least squares: Gauss-Newton and Levenberg-Marquardt}
\end{itemize}

\subsection{Full teaching explanation}\label{appendix-d8a7ddb675eb-full-teaching-explanation}

\begin{itemize}
\tightlist
\item
  \href{../02-Mathematics-for-Visual-Computing/main.pdf\#nameddest=02-foundations-ii-ch021-jacobians-curvature-and-local-approximation}{Phase 9 teaching section (Mathematics for Visual Computing, Ch. 5)}
\item
  \href{../01-Bits-Programs-and-Reliable-Systems/main.pdf\#nameddest=appendix-adf48f5dc5c1}{Visual Recall/09 - Calculus and Optimization - Visual Recall.md (Bits, Programs, and Reliable Systems, Reference notes: 09 - Calculus and Optimization - Visual Recall)} · visual recall
\end{itemize}

\subsection{Source coverage}\label{appendix-d8a7ddb675eb-source-coverage}

Supplied lesson parts: 822 (T04 · M050), 823 (T04 · M050), 824 (T04 · M050), 825 (T04 · M050).
